\part{Obligations, Degradation, and Recovery}
\label{part:III}

Part III represents obligations through time and introduces viability kernels, safe degradation, compensating continuation, and corrective capacity.

\begin{housing}Tenancy and repair duties become obligations, an unheated unit may not be let, and viability is studied under a maintenance budget.\end{housing}

\chapter{Obligations Through Time}\label{ch:11}
Represents contracts, promises, legal duties, maintenance commitments, and informal responsibilities as obligations with bearers, beneficiaries, triggers, deadlines, and discharge conditions. Explains why obligations cannot generally be inferred from a balance sheet or a snapshot of resources.
\scaffold

\chapter{Required Functions and Optional Capabilities}\label{ch:12}
Distinguishes functions that must continue from capabilities that can be reduced or suspended, and examines how these classifications are established, challenged, and revised. A public-service example shows that declaring a function optional is itself a consequential institutional decision.
\chspec{\Pref{P06}, \Pref{P20}.}
\scaffold

\chapter{Viability Under Scarcity}\label{ch:13}
Introduces viability kernels and constrained reachability, and distinguishes a desirable target from the set of states from which obligations can still be fulfilled. Uncertainty is treated separately for deterministic, robust, and probabilistic models. Comparative statics for added duties, capacity, and exposure are stated together.
\[
K=\bigl\{x:\exists\pi\ \text{such that all trajectories permitted by the model under }\pi\ \text{satisfy the continuation constraints}\bigr\}.
\]
\chspec{\Pref{P03}, \Pref{P04}, \Pref{P23}, \Pref{P25}.}
\scaffold

\chapter{Compensating Continuation}\label{ch:14}
Links suspension to restoration obligations, deadlines, and reserved resources. Distinguishes credible compensation from promises that depend on capabilities already being consumed. Worked cases are deferred building repairs, reduced hospital services, and temporary infrastructure closures.
\chspec{\Pref{P07}.}
\scaffold

\chapter{Corrective Capacity and Recoverability}\label{ch:15}
Defines the resources, permissions, knowledge, and time required to correct a course of action. Examines systems that remain operational while losing the ability to recover from foreseeable disturbances. A theorem gives conditions under which locally admissible steps leave the admissible region over a longer horizon.
\chspec{\Pref{P05}, \Pref{P14}.}
\scaffold
